\chapter{Conclusion and Future Work}
A further avenue for algorithmic refinement concerns the threshold parameter $s$ governing the size of each partial neighbourhood in the Aingworth approximation. In this project, two configurations were evaluated: the theoretically-derived default $s = \lfloor\sqrt{n \log n}\rfloor$ and a fixed override of $s = 10$. The latter was selected as a coarse empirical probe rather than as a principled tuning exercise. A natural extension would be to conduct a systematic parameter sweep --- evaluating, for instance, $s \in \{10, 25, 50, \lfloor\sqrt{n \log n}\rfloor / 2, \lfloor\sqrt{n \log n}\rfloor\}$ --- across each topological class and measuring the joint accuracy-runtime Pareto frontier. Such a study could identify topology-specific calibration heuristics, potentially yielding configurations that outperform the theoretical default on hub-dominated or grid-like topologies without sacrificing the formal approximation guarantee.


The project's conclusions should list the key things that have been learnt as a consequence of engaging in your project work. For example, ``The use of overloading in C++ provides a very elegant mechanism for transparent parallelisation of sequential programs'', or ``The overheads of linear-time n-body algorithms makes them computationally less efficient than $O(n \log n)$ algorithms for systems with less than 100000 particles''. Avoid tedious personal reflections like ``I learned a lot about C++ programming...'', or ``Simulating colliding galaxies can be real fun...''. It is common to finish the report by listing ways in which the project can be taken further. This might, for example, be a plan for turning a piece of software or hardware into a marketable product, or a set of ideas for possibly turning your project into an MPhil or PhD.